% Généré par tests/make_anesthesia_doc.py — ne pas éditer à la main.
\subsection{Agents IV, IM et locorégionaux}
\begin{tabularx}{\linewidth}{P{2.4cm}P{0.6cm}L}
\toprule
Agent & Niv. & Provenance (bloc \vp{anesthesia}) et notes \\
\midrule
Acépromazine & \evD & PK entièrement placeholder D (seule la t$\frac12$ IV de 7,1 h est publiée, Hashem 1992). Épargne de MAC plafonnée à 40 \% (halothane, Heard 1986, résumé : maximum dès 0,04 mg/kg) ; EC50 placeholder. Sédation placeholder. Blocage $\alpha$1 : veinodilatation (ven\_cap $\times$ (1 + 2$\cdot$E), borne du balayage) pour une PAM $-$12 \% chez le vigile (cible $-$17 \% sans changement de RVS, Rangel 2020 : écart signalé) ; hémoglobine $-$21 \% à effet maximal (CaO2 $-$21 \%, Rangel 2020 ; séquestration splénique, Sutil 2017). Pas d'analgésie. \\
Alfaxalone & \evC & PK : Vss 0,94 L/kg, Cl 26 mL/kg/min, Q 8,5 mL/kg/min (Dehuisser 2019, résumé) ; partage Vc/Vp placeholder D. ke0 placeholder D (classe). c\_imm calé pour que la dose d'induction publiée sans prémédication (2,6--2,8 mg/kg, Maney 2013, Hampton 2019) atteigne la perte de conscience avec la même marge que le propofol à 5,2 mg/kg (calage par session). Hémodynamique : PAM peu modifiée (Sarotti 2024, lu ; hypotension 1/14), FC +54 \% (Hampton 2019) : vasodilatation modérée placeholder D. Apnée : commande ventilatoire max(0, 1 $-$ 0,92$\cdot$EP) portée par le pic plasmatique, calée sur la dépendance à la vitesse d'injection (100 \% à 2 mg/kg/min contre 25 \% à 0,5 mg/kg/min, Bigby 2017) : hypopnée profonde en injection rapide après prémédication, ventilation conservée en injection lente. \\
Atipamézole & \evC & PK IM : Cl/F 554 mL/h/kg (9,23 mL/min/kg en supposant F = 1, D), Vz/F 1,49 L/kg (partage Vc/Vp D), ka 9,56 h$^{-1}$ calculé sur le Tmax de 0,35 h et la t$\frac12$ de 1,8 h (Cowan 2026, lu). Ki et ke0 placeholders D (éveil en 3 à 7 min après atipamézole IM, Vainio \& Vähä-Vahe 1990, résumé). Antagonisme compétitif de la dexmédétomidine (EC50 de toutes ses courbes multipliées par 1 + C/Ki). Rechute de sédation possible quand l'atipamézole s'élimine avant l'agoniste (Turunen 2019). \\
Bupivacaïne (bloc locorégional) & \evD & Module de bloc uniquement : délai, durée et efficacité placeholders D. PK systémique placeholder sans effet. À n'utiliser qu'en bloc locorégional (événement local\_block). \\
Dexmédétomidine & \evC & Clairance et Vss : Dent 2019 (résumé). Vc/Vp : partage placeholder D (Vc + Vp = Vss). IM : ka 5,96 h$^{-1}$ calculé sur le Tmax de 20,5 min (Aarnes 2023, résumé) avec t$\frac12$ 36 min ; F = 1 placeholder D. ke0 placeholder D (pic de sédation 12 min après IV, Alves 2024). Épargne de MAC (recalée le 20/11/2026) : EC50 3,0 ng/mL et pente 1,27 (max 0,8 imposé, D), ajustées sur les protocoles de Pascoe 2006 (charge puis CRI de 0,5 et 3 \textmu{}g/kg/h : $-$18 \% et $-$59 \%, cité par Acevedo-Arcique 2014, lu) aux concentrations que produit la PK du modèle à ces débits (1,1 et 6,7 ng/mL à 60 min). Les concentrations mesurées par Pascoe (0,198 et 1,903 ng/mL, résumé) sont 3 à 6 fois plus basses que celles du modèle : la PK (Dent 2019) n'est pas cohérente avec elles, écart non résolu (docs/VALIDATION\_HORS\_CALIBRAGE.tex). Hors échantillon : 2 \textmu{}g/kg + 2 \textmu{}g/kg/h, $-$43,4 \% (Acevedo-Arcique 2014, lu). Sédation : EC50 de l'ordre des concentrations au pic de sédation (Kuusela 2000 : 5,5 ng/mL à 10 \textmu{}g/kg) ; plafond 0,5 placeholder D (dose de propofol 5,2 $\to$ 2,9 mg/kg après 5 \textmu{}g/kg IM, Canfrán 2016). Antinociception : placeholder D. Effet biphasique : vasoconstriction périphérique ($\alpha$2B) portée par la concentration PLASMATIQUE (EP), sympatholyse centrale et tonus vagal par le site effet (E). Coefficients calés par petite grille (examples/anesthesia\_timing.rs) sur 10 \textmu{}g/kg IM à 15 min : modèle FC $-$58 \%, DC $-$57 \%, PAM +32 \% pour des cibles FC $-$55 \% et DC $-$59 \% (Congdon 2011) et PAM +36 \% (Alves 2024, 2 \textmu{}g/kg IV). Un bloc AV du 2e degré peut émerger du tonus vagal (av\_vagal du moteur) à plus forte dose, sans terme ad hoc ; il n'apparaît pas dans ce scénario de calage, alors que Karna 2022 le décrit chez tous les chiens (avec morphine). \\
Fentanyl & \evC & PK : Sano 2006 (résumé) : Cl 77,9 mL/min/kg, Vd 5 L/kg, t$\frac12$$\alpha$ 4,5 min, t$\frac12$$\beta$ 45,7 min ; Vc 4,05, Vp 0,95 L/kg et Q 116 mL/min/kg calculés (seule solution bicompartimentale compatible avec ces quatre valeurs). Épargne de MAC : plafond 84 \% et EC50 2,64 ng/mL (Machado 2022, résumé) ; pente 1 placeholder D. EC50 hémodynamique 3,65 ng/mL (FC, Machado 2022). Antinociception et sédation : EC50 reprises de la MAC, plafonds D ; ke0 D. Effets hémodynamiques et ventilatoires : forme et amplitudes du module opioide (classe, D pour le fentanyl), EC50 propre. \\
Kétamine & \evC & PK : Vc 0,371 L/kg, Vss 4,06 L/kg, Cl 58,2 mL/min/kg (Pypendop \& Ilkiw 2005, résumé) ; Q placeholder D. Épargne de MAC de l'isoflurane : Emax 92,2 \%, EC50 2,945 \textmu{}g/mL, $\gamma$ 3,01 (Pypendop 2007, résumé, sur les données de Solano 2006). c\_imm (perte de conscience) calé sur la dose d'induction de 6,1 mg/kg avec diazépam (Hampton 2019) : D. Antinociception, effet anti-hyperalgésique (NMDA) à faible concentration et ke0 : placeholders D. Anesthésie dissociative : réflexe palpébral conservé, globe central. Stimulation sympathique (s0 + 0,8$\cdot$E, plafond du balayage) : FC +23 \% après 6,1 mg/kg + midazolam pour une cible de +75 \% (Hampton 2019, kétamine-diazépam) ; le baroréflexe du moteur compense : écart signalé. Remarque : la contribution hypnotique (c\_imm) et l'épargne de MAC se cumulent dans la profondeur $\Delta$ : simplification signalée. \\
Lidocaïne & \evC & PK : Cl 27,5 mL/min/kg, Vss 1,38 L/kg (Ngo 1997, résumé) ; partage Vc/Vp et Q placeholders D. IM : ka 7,74 h$^{-1}$, fraction absorbée 91,9 \% (Wilcke 1983, résumé). Épargne de MAC : Emax ajusté sur Valverde 2004 (1,5 \textmu{}g/mL $-$18,7 \% ; 4,5 \textmu{}g/mL $-$43,3 \%, résumé) avec plafond 0,7 imposé (D) : EC50 3,15 mg/L, pente 1,36. Bloc locorégional (délai, durée, efficacité) et antinociception systémique : placeholders D. Toxicité non modélisée (le moteur n'a pas d'état convulsif) : extension tonique à 8,21 \textmu{}g/mL chez le chien (Wilcke 1983) ; ne pas dépasser cette concentration dans les scénarios. \\
Méthadone & \evC & SC : F 79 \% et ka 2,16 h$^{-1}$ calculé sur le Tmax de 1,26 h (Ingvast-Larsson 2010, résumé) ; IM non publiée : valeurs SC reprises (D). Épargne de MAC calée sur Credie 2010 (1 mg/kg : $-$48 \% à 2,5 h et $-$30 \% à 5 h ; 0,5 mg/kg : $-$35 \% et $-$15 \%, résumé, voie supposée IV). Antinociception : EC50 17 ng/mL (antinociception au-delà de 17 ng/mL, Amon 2021) ; plafond D. Sédation et ke0 : placeholders D. PK et PD hémodynamique reprises du module opioide 1.1.0 (méthadone, chien vigile). Sous halogéné, la dépression ventilatoire s'additionne à celle de l'agent inhalé (pas de terme de synergie). \\
Midazolam & \evC & PK : Vd 3,0 L/kg, Cl 27 mL/kg/min (Court \& Greenblatt 1992, résumé) ; partage Vc/Vp, Q, IM et ke0 placeholders D. Sédation plafonnée à 0,5 : la cible de propofol à l'intubation passe de 3,0 à 1,5 \textmu{}g/mL avec 0,2 mg/kg (Minghella 2016, résumé). Épargne de MAC : enflurane $-$55 \% à 949 ng/mL (Hall 1988, résumé) ; plafond 0,6 et EC50 D. Peu d'effets hémodynamiques ; légère dépression ventilatoire placeholder D. L'excitation paradoxale du chien vigile n'est pas modélisée. \\
Naloxone & \evC & t$\frac12$ IV 37 min et Cmax 18,8 ng/mL après 0,04 mg/kg (Wahler 2019, résumé) : Vc $\approx$ 2,1 L/kg et Cl $\approx$ 58 mL/min/kg calculés (approximation, C/D) ; Vp, IM, Ki et ke0 placeholders D (antagonisation de la méthadone en < 5 min après 0,1 mg/kg IM, Mills 2025, lu). Renarcotisation possible si la naloxone s'élimine avant l'opioïde (t$\frac12$ 37 min contre 3,9 h pour la méthadone). \\
Propofol & \evC & Cl 58,6 mL/kg/min et Vss 4,9 L/kg (Nolan 1993, résumé) ; partage Vc 0,8 / Vp 4,1 L/kg et Q placeholders D (Vc non publié). ke0 0,696 min$^{-1}$ (Lee 2009, résumé). c\_imm 6,2 mg/L : ajusté sur la réduction de MAC du sévoflurane (1,3 / 2,5 / 4,4 \textmu{}g/mL $\to$ $-$20,5 / $-$43 / $-$68 \%, Singsank-Coats 2015, résumé ; contribution Ce/c\_imm $\approx$ 0,16 MAC par \textmu{}g/mL) ; il place la perte de conscience (indice 0,4) vers 2,5 \textmu{}g/mL, cohérent avec Musk 2005 (65 \% d'inductions à 2,5 \textmu{}g/mL, 100 \% à 3,5) et la cible de 3 \textmu{}g/mL de Beths 2001. Vasodilatation (r\_sys) et émoussement du baroréflexe (baro\_gain, mécanisme retenu, forme D) calés ensemble pour PAM $-$31 \% au nadir et FC +62 \% (88 $\to$ 61 mmHg, 65 $\to$ 105/min, Cattai 2018, lu, après acépromazine + méthadone) ; la tachycardie vient du baroréflexe du moteur. Effets hémodynamiques portés par la concentration PLASMATIQUE (EP) : nadir de PAM à 55 s, avant l'équilibre du site effet (intubation à 140 s, Cattai 2018). Apnée d'induction : commande ventilatoire max(0, 1 $-$ 1,23$\cdot$EP), portée par le pic PLASMATIQUE : une injection rapide l'abolit, et le CO2 ne la relance qu'une fois la concentration redescendue. Calée sur la dépendance à la vitesse (apnée chez 100 \% des chiens à 4 mg/kg/min contre 25 \% à 1 mg/kg/min, Bigby 2017 ; 59 à 64 \% à 4 mg/kg, Murison 2001) : 5,2 mg/kg en 30 s sans prémédication donnent 56 s d'apnée, 1,8 mg/kg à 1 mg/kg/min après prémédication aucune. La vitesse d'injection se règle par la durée du bolus (30 s par défaut, iv\_push\_s du modèle). \\
\bottomrule
\end{tabularx}

\subsection{Halogénés}
\begin{tabularx}{\linewidth}{P{2.4cm}P{0.6cm}L}
\toprule
Agent & Niv. & Valeurs et provenance \\
\midrule
Isoflurane & \evC & MAC 1.3 \%, sang/gaz 1.4, tissu/sang 1.57 / 2.9 / 45.0. MAC 1,3 \% (1,28 à 1,39 \% : Steffey \& Howland 1977, Kazama 1988, Hofmeister 2008, résumés). Sang/gaz canin 1,40 (Soares 2012, lu). Cerveau/sang 1,57 (humain, Yasuda 1989) ; muscle/sang 2,9 et graisse/sang 45 : ordres de grandeur humains classiques, NON LUS (D). Effets en fonction de M = pression artérielle / MAC : Ees VG et VD $-$23,5 \% par MAC (Hoka 1995 : $-$23 \% à 1 MAC, $-$48 \% à 2 MAC) ; vasodilatation r\_sys = max(0,1 ; 1 $-$ c$\cdot$M\textasciicircum{}p) et émoussement du baroréflexe baro\_gain $\times$ (1 $-$ b$\cdot$M) calés par session sur le sévoflurane de Bernard 1990 (PAM $-$22 \% à 1,2 MAC, $-$38 \% à 2 MAC, DC $-$17 \% à 2 MAC ; isoflurane « similaire »). La dépression du baroréflexe par les halogénés est le mécanisme retenu (non chiffrée dans les sources lues : forme D) ; commande ventilatoire max(0, 1 $-$ a$\cdot$M\textasciicircum{}n) calée pour PaCO2 46 mmHg à 1 MAC en ventilation spontanée (Steffey 1993) et apnée à 2,6 MAC (Galloway 2004). Baroréflexe non déprimé par l'agent (placeholder : effet non sourcé). \\
Sévoflurane & \evC & MAC 2.2 \%, sang/gaz 0.664, tissu/sang 1.7 / 3.1 / 48.0. MAC 2,2 \% (2,10 à 2,36 \% : Kazama 1988, Hofmeister 2008, Yamashita 2012, Murahata 2018). Sang/gaz canin 0,664 (Soares 2012, lu). Cerveau/sang 1,70 (humain, Yasuda 1989) ; muscle/sang 3,1 et graisse/sang 48 : ordres de grandeur humains NON LUS (D). Ees $-$20 \% par MAC (Hoka 1995 : $-$16 \% à 1 MAC, $-$41 \% à 2 MAC) ; vasodilatation max(0,1 ; 1 $-$ c$\cdot$M\textasciicircum{}p) et émoussement du baroréflexe (1 $-$ b$\cdot$M, mécanisme retenu, forme D) calés sur Bernard 1990 (PAM $-$22 \% à 1,2 MAC, $-$38 \% à 2 MAC, DC $-$17 \% à 2 MAC) ; commande ventilatoire calée pour l'apnée à 3,45 MAC (Galloway 2004), avec la même PaCO2 à 1 MAC que l'isoflurane (hypothèse D, aucune valeur canine publiée trouvée). \\
\bottomrule
\end{tabularx}

\subsection{Constantes du modèle}
\begin{tabularx}{\linewidth}{P{3.6cm}P{1.2cm}P{0.6cm}L}
\toprule
Constante & Valeur & Niv. & Source \\
\midrule
\vp{co2_store_ml_mmhg_kg} & 1.0 & \evD & Placeholder : capacité de stockage du CO2 choisie pour une montée de PaCO2 de l'ordre de 4 à 5 mmHg/min en apnée (VCO2 4,6 mL/kg/min du patient). À sourcer sur une étude d'apnée chez le chien. Cohérence vérifiée : en apnée totale sous O2, PaCO2 de 96 mmHg à l'arrêt (DeBehnke 1995 : 93 $\pm$ 19 mmHg à 11,4 min, résumé). \\
\vp{frc_ml_kg} & 45.0 & \evC & CRF 45 $\pm$ 8 mL/kg chez le chien sous halothane, curarisé, ventilé (Larsson 1992, Acta Anaesthesiol Scand 36:146, résumé) ; 30 à 75 mL/kg selon la posture et la sédation chez le chien vigile (Rozanski 2010, PMC2801309, lu). \\
\vp{atelectasis_aa_mmhg} & 110.0 & \evC & Gradient A-a ajouté sous anesthésie, proportionnel à la profondeur et à (FiO2 $-$ 0,21)/0,79 : effet du shunt d'atélectasie, négligeable à l'air. Calé le 20/11/2026 : PaO2 547 [495--616] mmHg sous isoflurane 1,3--1,5 \% en ventilation contrôlée et O2 (Itami 2019, doi:10.1371/journal.pone.0213414, lu ; modèle 548) ; 523 (513--539) mmHg sous sévoflurane 1,3 MAC en ventilation spontanée (Kato 2024, doi:10.5455/OVJ.2024.v14.i5.20, lu ; modèle 523) ; 570--579 chez Mosing 2025 (lu). Avant : 630 à 660 mmHg. \\
\vp{vd_ml_kg} & 3.0 & \evD & Placeholder : espace mort anatomique + appareil en ventilation contrôlée. \\
\vp{circuit_l} & 3.0 & \evD & Placeholder : volume d'un circuit filtre adulte. \\
\vp{etco2_gap_mmhg} & 3.0 & \evD & Placeholder : écart PaCO2 $-$ EtCO2 chez le sujet sain. \\
\vp{flowby_fio2} & 0.4 & \evD & Placeholder : FiO2 obtenue en flow-by. \\
\vp{loc_mac} & 0.65 & \evC & MAC-awake du chien : rapports 0,77 (isoflurane 1,0/1,3 \%) et 0,62 (sévoflurane 1,3/2,1 \%), Hofmeister 2008 (résumé) ; moyenne arrondie. \\
\vp{awareness_hill} & 3.0 & \evC & Pente de la conscience sous la MAC-awake, awareness = 1 $-$ (indice/loc\_mac)\textasciicircum{}n. Calée le 20/11/2026 sur les délais de réveil (décubitus sternal) : isoflurane FE 1,3 \% ($\approx$ 1,0 MAC) 2 h, sternal 14,7--18,0 min (Rastas 2025, PMC11781748, lu ; modèle 17,2) ; sévoflurane $\approx$ 105--120 min, extubation 7 (5--7) et sternal 9,5 (7,25--11,75) min (Lozano 2009, résumé ; modèle à 1,45 MAC : 9,0 et 12,1). Avec n = 1 (avant) : sternal à 104 min. \\
\vp{swallow_mac} & 0.5 & \evC & MAC-extubation du sévoflurane : rapport 0,51 $\pm$ 0,08 (le chien lève la tête ou mâchonne la sonde), Murahata 2018 (PMC6068302, lu). \\
\vp{palpebral_mac} & 0.9 & \evD & Placeholder : perte du réflexe palpébral. \\
\vp{jaw_mac} & 1.2 & \evD & Placeholder : perte du tonus de la mâchoire. \\
\vp{surgical_mac} & 1.0 & \evD & Placeholder : début du plan chirurgical (1 MAC : 50 \% d'immobilité au stimulus supramaximal). \\
\vp{deep_mac} & 1.8 & \evD & Placeholder : plan profond (globe recentré). \\
\vp{excessive_mac} & 2.5 & \evD & Placeholder : surdosage (mydriase). \\
\vp{mac_bar} & 1.6 & \evC & MAC-BAR du sévoflurane : 3,33 \% pour une MAC de 2,10 \% (rapport 1,59, Yamashita 2012, résumé) ; rapport 2,05 chez Murahata 2018. \\
\vp{bar_hill} & 3.0 & \evD & Placeholder : pente de la suppression de la réponse autonome. \\
\vp{move_hill} & 10.0 & \evD & Placeholder : pente de la courbe d'immobilité (relation dose-réponse raide de la MAC). \\
\vp{tau_nu_s} & 15.0 & \evD & Placeholder : constante de temps de la réponse nociceptive. \\
\vp{tau_stress_min} & 20.0 & \evD & Placeholder : constante de temps de la réponse de stress (VO2, glycémie). \\
\vp{sens_rate_per_min} & 0.05 & \evD & Placeholder : vitesse de sensibilisation centrale sous nociception non contrôlée. \\
\vp{sens_tau_h} & 24.0 & \evD & Placeholder : retour de la sensibilisation. \\
\vp{sens_gain} & 0.5 & \evD & Placeholder : amplification maximale de la nociception par la sensibilisation (+50 \%). \\
\vp{t0_c} & 38.5 & \evD & Placeholder : température de référence du chien. \\
\vp{t_amb_c} & 21.0 & \evD & Placeholder : température du bloc opératoire. \\
\vp{tau_cool_h} & 12.0 & \evD & Placeholder : constante thermique ; avec loss\_gain\_anesth = 2 et VO2 $-$20 \%, refroidissement initial $\approx$ 1,75 °C/h. \\
\vp{loss_gain_anesth} & 3.5 & \evC & Pertes thermiques multipliées sous anesthésie (vasodilatation, perte du comportement), au-dessus du seuil de vasoconstriction. Calé le 20/11/2026 avec vasoconstriction\_c sur Redondo 2012 (1525 chiens, résumé) : $-$2,0, $-$2,6 et $-$2,9 °C à 60, 120 et 180 min (modèle $-$2,0, $-$2,6, $-$2,8 ; isoflurane 1,3 MAC, salle 21 °C) ; chute rapide puis plateau après 60 min (Sage 2023, 1097 chiens, résumé). Avant : 2,0 sans plateau ($-$4 °C à 3 h). \\
\vp{vasoconstriction_c} & 35.5 & \evC & Température où l'excès de pertes de l'anesthésie s'éteint (vasoconstriction thermorégulatrice) : décroissance linéaire de l'excès entre t0\_c et ce seuil. Calé le 20/11/2026 avec loss\_gain\_anesth sur le plateau de Redondo 2012 (35,8 $\pm$ 1,5 °C à 180 min, résumé) et Sage 2023 (résumé). Forme D. \\
\vp{warming_c_per_h} & 3.5 & \evC & Apport d'un réchauffement actif (matelas à eau + air pulsé). Calé le 20/11/2026 : température finale 38,2 $\pm$ 1,0 et 38,7 $\pm$ 0,5 °C avec réchauffement contre 35,4--35,7 sans, après $\approx$ 2 h d'isoflurane (Rastas 2025, PMC11781748, lu ; modèle 38,4 à 3 h). Avant : 1,5 °C/h (insuffisant avec les pertes recalées). \\
\vp{mac_per_c} & 0.05 & \evD & Placeholder : réduction de MAC par °C d'hypothermie (ordre de grandeur humain non lu). \\
\vp{blood_volume_ml_kg} & 85.0 & \evD & Volume sanguin de 85 mL/kg, déjà utilisé par le baroréflexe du moteur (MODEL\_SPEC \S{}3). \\
\vp{crystalloid_tau_min} & 30.0 & \evD & Placeholder : sortie vasculaire des cristalloïdes. \\
\vp{refill_frac} & 0.6 & \evC & Fraction du volume perdu restituée au plasma par remplissage transcapillaire ; avec refill\_tau\_h, 30 \% à 2 h (Joles 1989, chiens conscients, résumé) et 56 \% à 8 h (Miller \& Dale 1978, résumé). \\
\vp{refill_tau_h} & 2.9 & \evC & Constante du remplissage transcapillaire, ajustée avec refill\_frac sur Joles 1989 et Miller \& Dale 1978 (résumés). \\
\vp{spleen_ml_kg} & 6.8 & \evC & Sang mobilisable par la rate : 8 \% du volume sanguin en hémorragie chez le chien conscient (Guntheroth \& Mullins 1963, résumé) ; 2,4--3,2 mL/kg dans les réflexes et hémorragies modérées (Shoukas 1981, Carneiro \& Donald 1977). \\
\vp{spleen_tau_min} & 1.0 & \evD & Placeholder : contraction splénique en une à quelques minutes. \\
\vp{spleen_s_span} & 0.3 & \evD & Placeholder : hausse du tonus sympathique qui vide entièrement la rate. \\
\vp{spleen_hct_ratio} & 1.8 & \evD & Placeholder : sang splénique concentré (Ht 28,9 \% chez les splénectomisés contre 39,8 \% chez les témoins après 50 \% de saignée, Hoekstra 1988, résumé). \\
\vp{pressor_gain} & 0.6 & \evC & Hausse maximale des RVS par vasopressine et angiotensine II en hémorragie. Calée le 21/11/2026 sur le chien conscient : PAM $-$10 mmHg à 17 mL/kg (Shen 1990 : $-$15 $\pm$ 3 à 20 mL/kg), $-$23 mmHg à 30 mL/kg (Thrasher 1998 : $-$35 $\pm$ 10), DC $-$41 \% à 25 \% (Francis 2004 : $-$40 \%, résumés). Le blocage de l'angiotensine aggrave la chute (losartan : $-$40 \% au lieu de $-$25 \%, Francis 2004 ; saralasine : PAM 102 $\to$ 80, Spielman 1975). \\
\vp{pressor_lost_frac} & 0.25 & \evC & Mi-effet presseur vers 25 \% du volume sanguin : vasopressine 6, 200 puis 991 pg/mL à 10, 20 et 30 mL/kg (Wang 1988, résumé). \\
\vp{pressor_hill} & 4.0 & \evD & Placeholder : montée raide de la vasopressine entre 10 et 30 mL/kg. \\
\vp{iv_push_s} & 30.0 & \evC & Durée d'injection par défaut d'un bolus IV : 30 s, comme l'induction au propofol de Cattai 2018 (PMC6134702, lu). La vitesse d'injection change l'incidence de l'apnée (Murison 2001, Bigby 2017) et l'hypotension : elle est réglable à chaque injection. \\
\vp{o2_debt_tol_ml_kg} & 20.0 & \evD & Placeholder : dette d'O2 tolérée avant l'arrêt hypoxique (réserves, anaérobiose), environ 4 min de VO2 totale (5,5 mL/kg/min) : l'arrêt survient alors quelques minutes après le début d'une apnée à l'air ambiant. Le moteur stationnaire, lui, conclut au décès dès que VO2 dépasse DO2. Délai d'arrêt du modèle en apnée totale sous O2 : 10,3 min depuis l'atélectasie du 20/11/2026 (11,6 avant ; DeBehnke 1995 : 11,4 $\pm$ 2,4 min, chiens anesthésiés, sonde clampée). \\
\vp{o2_debt_tau_min} & 10.0 & \evD & Placeholder : constante de remboursement de la dette d'O2 une fois l'apport rétabli. \\
\vp{shock_tol_min} & 60.0 & \evC & Durée cumulée sous map\_arrest (PAM < 40 mmHg) avant le choc irréversible : 1 à 3 h chez le chien maintenu vers 40 mmHg (Morales 1999, Granot 1985, résumés, déjà cités pour map\_arrest), borne basse retenue. Le moteur stationnaire conclut au décès dès PAM < 40 ; la session tolère cette hypotension pendant ce délai. Chiens survivant à PAM 40--44 mmHg et DC 0,6 L/min sous isoflurane (Cardillo 2024, PMC11551537, lu). Ajouté le 20/11/2026 (docs/VALIDATION\_HORS\_CALIBRAGE.tex). \\
\vp{do2_crit_ml_kg_min} & 9.8 & \evC & Apport d'O2 critique 9,8 mL/kg/min chez le chien anesthésié, hypoxie hypoxique et anémique (Cain 1977, J Appl Physiol 42:228, résumé) ; 7,7 mL/kg/min en hémorragie (Schertel 1994, résumé). \\
\vp{blood_o2_store_frac} & 0.5 & \evD & Fraction du sang dont le contenu en O2 suit l'alvéole pendant l'apnée (réserve sanguine, pente de la courbe de dissociation) : placeholder calé sur le délai d'arrêt par asphyxie (modèle 11,6 min sous O2, DeBehnke 1995 : 11,4 $\pm$ 2,4 min, résumé) ; sans effet à l'équilibre. \\
\bottomrule
\end{tabularx}

\subsection{Scénarios chirurgicaux}
\textbf{Orchiectomie (chien)} (\evD) : Incision prescrotale (0.8, 60 s) ; Ligature du premier cordon spermatique (1.1, 120 s) ; Ligature du second cordon spermatique (1.1, 120 s) ; Fermeture (0.5, 300 s). Douleur postopératoire 0.4, constante 8 h. Intensités relatives : placeholders pédagogiques D. Le bloc intratesticulaire se pose sur le territoire « spermatic\_cord ».

\textbf{Ovariohystérectomie (chienne)} (\evD) : Incision cutanée et de la ligne blanche (0.8, 60 s) ; Extériorisation de la corne utérine (0.7, 120 s) ; Rupture du ligament suspenseur et ligature du pédicule ovarien (1.2, 180 s) ; Second pédicule ovarien (1.2, 180 s) ; Ligature du col utérin (0.9, 120 s) ; Fermeture de la paroi et de la peau (0.5, 600 s). Douleur postopératoire 0.6, constante 12 h. Intensités relatives (1 = stimulus supramaximal de la définition de la MAC) : placeholders pédagogiques D qui ordonnent les temps opératoires (le pédicule ovarien est le temps le plus stimulant, enseignement classique). Aucune valeur publiée.

